\section{Pruebas y rendimiento}
\label{sec:pruebas}
\textit{Responsables: todo el equipo (QA-01); Mariano (PERF-01); Fabricio (PERF-02).}

\subsection{Pruebas automatizadas}
Cada microservicio tiene pruebas con moto (AWS simulado) que cubren permisos por rol, aislamiento por
tenant, validación del cuerpo, idempotencia del consumidor y el recorrido del evento
(outbox $\rightarrow$ Stream $\rightarrow$ bus $\rightarrow$ regla $\rightarrow$ cola). La plantilla de
servicio trae 15 pruebas base. También se puede probar en local con LocalStack (LOC-01).

\pendiente{Cada dueño: número de pruebas de su microservicio. Nota: a fines de noviembre algunas
pruebas fallan en CI por la versión nueva de moto (\texttt{DescribeStream: Not yet implemented});
indicar si se fijó la versión.}

\subsection{Prueba de punta a punta del curso piloto (QA-01)}
\pendiente{Resultado del recorrido con 2 tenants (UTEC y UPC): contenido $\rightarrow$ tarea
$\rightarrow$ entrega $\rightarrow$ nota $\rightarrow$ aviso. Ver \texttt{qa/piloto.py} en
\texttt{nubeduca-backend}.}

\subsection{Prueba de carga (PERF-01)}
\texttt{perf/k6/carga\_cursos.js} (k6) lista cursos y abre uno, repartiendo las peticiones entre
UTEC, UPC y PUCP en proporción 15/10/5 y usando como usuarios a los docentes del seed. Tiene dos
perfiles: \texttt{corto} (5 usuarios virtuales, 1{,}5\,min, para \texttt{dev}) y \texttt{estres}
(hasta 100 usuarios virtuales, para staging). \texttt{perf/tokens/generar\_tokens.sh} obtiene ID tokens reales para el modo Cognito.

Corrida corta del 20-Nov-2026 contra ms-cursos en \texttt{dev} con 30\,000 cursos:

\begin{table}[htbp]
\centering\small
\begin{tabularx}{\textwidth}{Xll}
\toprule
Métrica & Resultado & Meta \\
\midrule
Peticiones & 273 (unas 3 por segundo) & --- \\
Errores & 0\,\% & $<$ 1\,\% \\
\texttt{GET /courses/\{courseId\}}, p95 & 43\,ms & $<$ 1\,s \\
\texttt{GET /courses}, p95 & 1{,}84\,s & $<$ 1\,s (no se cumple) \\
\bottomrule
\end{tabularx}
\caption{Resultados de PERF-01. Las metas son provisionales.}
\label{tab:perf01}
\end{table}

El listado tarda unos 1{,}7\,s incluso con un solo usuario, así que no es efecto de la carga sino de
la consulta: filtra \texttt{ownerId} sobre toda la partición del tenant sin índice
(sección~\ref{sec:ms-cursos}). Con un índice por docente se espera un p95 de 50 a 100\,ms.

\evidencia{perf01-reporte.png}{Reporte de k6 de la corrida corta de PERF-01.}{fig:perf01}

\pendiente{Ajustar las metas si la sección 13 del documento de arquitectura define otras.}

\subsection{Prueba de estrés en staging (PERF-02)}
\pendiente{Fabricio: resultados de \texttt{PERFIL=estres} en la cuenta de integración (opcional).}
